\titlespacing*{\subsection}{0pt}{1.0em}{0.3em}
\titlespacing*{\paragraph}{0pt}{1.0em}{0.3em}

본 연구는 전통적 시계열 기반 환율 예측 모형이 충분히 반영하지 못했던 국가 간 상호의존 구조를 체계적으로 포착하고, 환율 변동의 복합적인 변동 구조를 정교하게 모델링하는 데 초점을 둔다. 이를 위해 베이지안 다변량 시계열 그래프 신경망을 기반으로 IMF 주요 3개국(미국, 일본, 한국)의 월간 환율 데이터와 거시경제 지표를 결합하여 시공간적 상관 구조를 통합적으로 분석하였다. 특히 다변량 환율 및 거시경제 시계열에 내재된 잠재적 상호연결성을 학습하기 위해 적응형 그래프 학습 모듈을 도입하였으며, 이를 통해 예측 성능의 개선과 모델의 확장 가능성을 실증적으로 확인하였다. 더 나아가, 단순 수치 예측을 넘어 경제적 해석을 가능하게 하기 위해 멀티 에이전트 시스템을 통합하였다. 이를 통해 예측된 환율과 거시경제 요인 간의 구조적 관계를 해석하고, 경제적 맥락에 기반한 논리적 설명을 제시함으로써 인공지능 기반 모델의 설명 가능성을 강화하였다.
\\[0.5em]
\textbf{주제어}: 베이지안 다변량 시계열 그래프, 멀티 에이전트 시스템, 설명 가능한 인공지능, 환율 변동 구조 모델링
\\[0.8em]
\textbf{Abstract.} This study focuses on systematically capturing the interdependence structures across nations, which traditional time-series based exchange rate forecasting models have struggled to fully incorporate and elaborately modeling the complex dynamics of exchange rate fluctuations. To this end, we propose a framework based on a Bayesian Multivariate Time-series Graph Neural Network (B-MTGNN) to comprehensively analyze spatiotemporal correlation structures by integrating monthly exchange rate data and macroeconomic indicators from three major IMF member countries. In particular, an adaptive graph module is introduced to learn the latent interconnections inherent in multivariate exchange rate and macroeconomic time series, through which we empirically demonstrate improvements in predictive performance and model scalability. Furthermore, to move beyond simple numerical forecasting and enable economic interpretation, we propose integrating a Multi-Agent System (MAS). This approach enhances the explainability of the AI-based model by interpreting the structural relationships between predicted exchange rates and macroeconomic factors, thereby providing logical explanations grounded in an economic context.
\\[0.5em]
\textbf{Keywords}: Bayesian Multivariate Time-series Graph Neural Network (B-MTGNN), Multi-Agent System (MAS), Explainable AI (XAI), Exchange Rate Dynamics Modeling
